\documentclass[a4paper]{article}
\usepackage{amsfonts}
\usepackage{amsmath}
\usepackage{amssymb}
\usepackage{graphicx}     

\begin{document}
  
\noindent \large Auteur: Noran Ben Said \\
\noindent \large Studierichting: Informatica\\
\noindent \large Oefeningenreeks 6B nr. 17.7\\

\medskip

\normalsize


$1^3 + 3^3 + 5^3 + \dots + (2n-1)^3$\\

$s_n = \displaystyle\sum_{k=1}^{n} (2k-1)^3$\\

$(2k-1)^3 = (2k)^3 - 3(2k)^2(1) + 3(2k)(1)^2 - 1^3$\\

$= 8k^3 - 12k^2 + 6k - 1$\\

$s_n = \displaystyle\sum_{k=1}^{n} (8k^3 - 12k^2 + 6k - 1)$\\

$= 8\displaystyle\sum_{k=1}^{n} k^3 - 12\displaystyle\sum_{k=1}^{n} k^2 + 6\displaystyle\sum_{k=1}^{n} k - \displaystyle\sum_{k=1}^{n} 1$\\

$\displaystyle\sum_{k=1}^{n} k^3 = \left(\dfrac{n(n+1)}{2}\right)^2 = \dfrac{n^2(n+1)^2}{4}$\\

$\displaystyle\sum_{k=1}^{n} k^2 = \dfrac{n(n+1)(2n+1)}{6}$\\

$\displaystyle\sum_{k=1}^{n} k = \dfrac{n(n+1)}{2}$\\

$\displaystyle\sum_{k=1}^{n} 1 = n$\\


$s_n = 8 \cdot \dfrac{n^2(n+1)^2}{4} - 12 \cdot \dfrac{n(n+1)(2n+1)}{6} + 6 \cdot \dfrac{n(n+1)}{2} - n$\\

$= 2n^2(n+1)^2 - 2n(n+1)(2n+1) + 3n(n+1) - n$\\


$= 2n^2(n^2 + 2n + 1) - 2n(n+1)(2n+1) + 3n(n+1) - n$\\

$= 2n^4 + 4n^3 + 2n^2 - 2n(2n^2 + 3n + 1) + 3n^2 + 3n - n$\\

$= 2n^4 + 4n^3 + 2n^2 - 4n^3 - 6n^2 - 2n + 3n^2 + 2n$\\

$= 2n^4 - n^2$\\

$= n^2(2n^2 - 1)$\\


$s_n = n^2(2n^2 - 1)$ voor alle positieve gehele getallen
n.\\

\begin{thebibliography}{999}
%\bibitem{a} Referentie
\end{thebibliography}
\end{document} 
